\begin{textsample}{POS Dim 3 – rs – Score 20.00 – rs/rs\_inf\_7.txt}  \label{ex:f3_pos_245}
7.
OS PROCESSOS DE ACOLHIMENTO E TRANSIÇÃO DE \textbf{BEBÊS}, \textbf{CRIANÇAS} BEM \textbf{PEQUENAS} E \textbf{CRIANÇAS} \textbf{PEQUENAS} Na Educação \textbf{Infantil}, há diferentes transições a serem consideradas, como os processos transitórios de casa para \textbf{instituição} de Educação \textbf{Infantil}, transições no interior da \textbf{instituição}, transição Creche/Pré-escola e transição Pré-escola/Ensino Fundamental.
Considerar o ingresso de \textbf{bebês} e \textbf{crianças} na Educação \textbf{Infantil} demanda um processo de acolhimento e adaptação que envolve as próprias \textbf{crianças}, as famílias e os profissionais da escola.
Assim como deve prever ações de transição e acolhimento para a entrada das \textbf{crianças} no Ensino Fundamental.
A adaptação ocorre sempre que a \textbf{criança} se \textbf{depara} com uma nova etapa ou um novo ambiente educativo, podendo ser em relação à mudança de escola, de turma, de professor referência ou mesmo entre os diferentes momentos da jornada \textbf{diária}.
Na Educação \textbf{Infantil}, o novo por vezes gera insegurança, visto que os \textbf{bebês} e as \textbf{crianças}, que vivem exclusivamente em seus contextos familiares, \textbf{deparam} -se com a diversidade presente em um ambiente coletivo, com um funcionamento diferente do habitual em seus lares, passando a participar de atividades incomuns ao seu cotidiano e a conviver com \textbf{adultos} e \textbf{crianças} inicialmente estranhos.
Entende -se, então, que a adaptação é considerada como um momento de transição, tendo em vista que, de maneira gradativa, a \textbf{criança} vai criando vínculos com professores e outros \textbf{adultos}, com outras \textbf{crianças} e com o meio.
Esse período demanda sensibilidade e olhar \textbf{atento} do professor e demais profissionais da \textbf{instituição}, de modo que as necessidades das \textbf{crianças} sejam atendidas.
A adaptação na Educação \textbf{Infantil} precisa ser compreendida na perspectiva do acolhimento como princípio norteador para o trabalho educativo.
A organização do ambiente precisa ser pensada para \textbf{acolher} e motivar as aprendizagens das \textbf{crianças} ; as \textbf{rotinas} e as jornadas \textbf{diárias} precisam \textbf{acolher} as experiências dos \textbf{bebês} e das \textbf{crianças}, dando -lhes o tempo necessário para \textbf{brincar} e explorar ; o período de adaptação precisa \textbf{acolher} as \textbf{crianças} e suas famílias e levar em conta as emoções que surgem neste período e depois dele.
Não há contradição entre acolhimento e ação educativa, pois, \textbf{acolher} significa valorizar as manifestações das \textbf{crianças} e reconhecê -las como experiências reais, capazes de levá -las à construção de aprendizagens e de vínculos significativos.
Um dos elementos facilitadores para esse processo de adaptação são as propostas educativas, que, geralmente, priorizam uma organização peculiar dos espaços, dos tempos e da própria \textbf{rotina} cotidiana, ou seja, do ambiente compreendido numa dimensão que abrange tanto os espaços físicos quanto as relações que neles se estabelecem.
Assim compreendidos, os ambientes precisam ser planejados para \textbf{acolher} as atividades \textbf{lúdicas}, para oportunizar que as \textbf{crianças} realizem ações com autonomia, fazendo surgir situações \textbf{interessantes}, relações que permitem bem-estar, contextos que promovam a riqueza da \textbf{brincadeira} e a construção de vínculos entre as \textbf{crianças} e o professor.
O processo de transição das \textbf{crianças} que passam da Educação \textbf{Infantil} para o Ensino Fundamental não se restringe a uma simples transferência de ritos e propostas da Educação \textbf{Infantil} para o 1º ano do Ensino Fundamental, pois existem especificidades a serem consideradas em cada etapa.
É necessário refletir sobre um currículo voltado para a integração e a continuidade dos processos de aprendizagem das \textbf{crianças} entre as etapas, tendo como ênfase o acolhimento afetivo e a continuidade das aprendizagens realizadas pelas \textbf{crianças} na Educação \textbf{Infantil}, sem adotar práticas preparatórias ou antecipar processos de aprendizagem específicos da etapa seguinte, mas garantir as especificidades de cada momento do percurso educativo das \textbf{crianças}.
As DCNEB reafirmam que : > Art. 11.
Na transição para o Ensino Fundamental a proposta pedagógica deve prever formas para garantir a continuidade no processo de aprendizagem e desenvolvimento das \textbf{crianças}, respeitando as especificidades \textbf{etárias}, sem antecipação de conteúdos que serão trabalhados no Ensino Fundamental. ( DCNEB, 2013, p. 100 ).
Ao planejar a transição das \textbf{crianças} para o Ensino Fundamental, é preciso conhecer a trajetória educativa realizada na Educação \textbf{Infantil}, buscando informações contidas em relatórios, portfólios ou outros \textbf{registros} que documentem e evidenciem o desenvolvimento e as aprendizagens das \textbf{crianças}, bem como organizar visitas, compartilhar materiais e diálogos entre os professores das duas etapas, buscando evitar a fragmentação e a descontinuidade do trabalho pedagógico.
Envolver diretamente as \textbf{crianças} e o contexto familiar em ações de transição como visitas, reuniões, conversas, entrevistas, entre outras, são estratégias de acolhimento e adaptação que \textbf{apoiarão} as \textbf{crianças} a superarem os desafios da transição.
A BNCC afirma a necessidade de que as propostas pedagógicas garantam os direitos das \textbf{crianças}, tanto na Educação \textbf{Infantil} quanto no Ensino Fundamental.
A \textbf{infância} não se encerra quando as \textbf{crianças} completam 6 anos e ingressam na etapa seguinte, mas continua ao longo dos primeiros anos do Ensino Fundamental.
Considerando os Direitos e os Objetivos de aprendizagem e desenvolvimento, tanto a BNCC quanto o Referencial Curricular Gaúcho apresentam uma síntese das aprendizagens que são esperadas em cada Campo de Experiências, ao longo de toda a Educação \textbf{Infantil}, como um elemento balizador e \textbf{indicativo} para o planejamento das ações educativas desenvolvidas na escola da \textbf{infância}, a serem ampliadas e aprofundadas no Ensino Fundamental, jamais sendo utilizada como requisitos para o acesso ao Ensino Fundamental.

% matched lemmas: acolher, adulto, apoiar, atento, bebê, brincadeira, brincar, criança, deparar, diário, etário, indicativo, infantil, infância, instituição, interessante, lúdico, pequeno, registro, rotina
\end{textsample}
